%%%%%%%%%%%%%%%%%%%%%%%%%%%%%%%%%
%
%       EXERCÍCIO
%
%%%%%%%%%%%%%%%%%%%%%%%%%%%%%%%%%

\ifdefstring{\atividade}{prova}{%
    \ifdefstring{\modo}{objetivo}{%
        \renewcommand{\valorquestao}{\ValorQObj\ ponto}
    }{%
        \renewcommand{\valorquestao}{\ValorQDisc\ pontos}
    }%
}%

\begin{exercicioBanco}[\valorquestao]
Para organizar um grande estacionamento em um parque, um administrador precisa estimar a quantidade de funcionários necessários para controle de entrada e saída. Para cada 250 veículos estacionados, é necessário um funcionário.

Às 8h da manhã, o administrador observa que a área ocupada pelos veículos forma um retângulo de dimensões \(300\) m por \(200\) m. Sabe-se que, em média, cada veículo ocupa uma área de \(12\ \text{m}^2\).

Nas horas seguintes, espera-se que o número de veículos aumente a uma taxa de \(5.000\) veículos por hora até as 12h, quando o estacionamento será fechado.

Determine quantos funcionários serão necessários nesse horário.

% Define as alternativas
\newcommand{\alternativas}{%
    \begin{center}
        \begin{tabularx}{\textwidth}{XXXXX}
            (a) \(80\). &
            (b) \(100\). &
            (c) \(120\). &
            (d) \(140\). &
            (e) \(160\).
        \end{tabularx}
    \end{center}
}

% Resposta correta
\newcommand{\resposta}{B}

% Lógica condicional para exibição
\ifdefstring{\atividade}{lista}{%
    \alternativas
    \vspace{0.5em}
    
    \noindent\textbf{Resposta:} letra \textbf{\resposta}.
}{%
    \ifdefstring{\modo}{objetivo}{%
        \alternativas
    }{%
        % modo = discursiva → não mostra alternativas
    }
}
\end{exercicioBanco}

